% page 404 du fac-simile (image p-403 du scan)
% bloc 154.9 x 246.3 mm ; marge gauche 8.0 mm
\pgc{-12.700mm}{0.000mm}{
\l{\cel{3}396}
\l{}
\l{\sou{noktuo}. (\sou{zool.})\cel{2}Speco di strigo, rapt-ucelo noktala.}
\l{\cel{3}- L.\cel{1}\sou{noctua}.}
\l{}
\l{\sou{nokturno}. I. (\sou{muziko}). Sorti di romanco du-voca, kun la}
\l{\cel{3}karaktero di serenado. - II. (\sou{liturgio katol.}) Parto}
\l{\cel{3}di la nokt-oficio, qua kontenas ula quanto de psalmi.}
\l{\cel{3}-DEFIRS.}
\l{}
\l{\sou{nomar}.\cel{2}(\sou{trans.}) Dicar la nomo qua distingas (persono) ed}
\l{\cel{3}indikas lu individuale. - DEFIS.}
\l{}
\l{\sou{nomado}. Qua ne havas rezideyo fixa en teritorio determini-}
\l{\cel{3}ta. - DEFIRS.}
\l{}
\l{\sou{nombro}. (\sou{matem.}) Vorto qua indikas quanta-foye la unajo}
\l{\cel{3}kontenasas en quanto. - EFIS.}
\l{}
\l{\sou{nomenklaturo}. Ensemblo ek la termini, ek la nomi specala,}
\l{\cel{3}qui uzesas da ca o ta cienco, arto, libro teknikala.}
\l{\cel{3}- DEFIS.}
\l{}
\l{\sou{nominar}. (\sou{trans}.)Grantar titulo,ofico. - EFIS.}
\l{}
\l{\sou{nominativo}.\cel{2}(\sou{gram.)}\cel{2}I. Deklino-kazo qua indikas, ke la}
\l{\cel{3}nomo esas la subjekto di la verbo. - II. La ipsa nomo}
\l{\cel{3}qua esas la subjekto di la verbo. - DEFIS.}
\l{}
\l{\sou{nomografio}. (\sou{matem.}) Ensemblo ek la metodi qui havas kom}
\l{\cel{3}skopo remplasar la kalkuli numerala per simpla lekto}
\l{\cel{3}sur tabelo grafika di qua la nomo esas \textquotedbl{}abako\textquotedbl{}. -}
\l{}
\l{\sou{non}. (\sou{matem.})\cel{2}La cifro \textquotedbl{}9\textquotedbl{}.}
\l{}
\l{\sou{nona-}. - Prefixo ciencala = non-....}
\l{\cel{23}54}
\l{\sou{noniliono}. (\sou{matem.})\cel{2}l0\cel{2}, od oktiliono x l.000.000.}
\l{}
\l{\sou{nopalo.}\cel{2}(\sou{bot.)}\cel{3}Arboro sur qua vivas la kochenilo, varie-}
\l{\cel{3}tato di kaktuso. - L.\cel{1}\sou{opuntia cocoinillifera}. DEFIRS.}
\l{}
\l{\sou{nordo.}\cel{2}I. Un ek la quar punti kardinala, punto di horizonto}
\l{\cel{3}situita en la direciono di la polo-stelo. - II. Parto}
\l{\cel{3}di la Terglobo, regiono situita vicine a nordo. - DEFIRS.}
\l{}
\l{\sou{norio}.\cel{2}(\sou{teknol.})\cel{2}Mashino hidraulikala, qua konsistas ye}
\l{\cel{3}siteli ligita alonge kateno senfina, qua spulas sur tam-}
\l{\cel{3}buro, per rotaco, e per qua la siteli plneskas ed acensas,}
\l{\cel{3}ca-latere - vakueskas e decensas, ta-latere. - DEF.}
\l{}
\l{\sou{normo}. I. Regulo, principo di konduto. - II. (\sou{industrio}).}
\l{\cel{3}Regulo qua fixigas la cirkonstanci di la faco di objekto}
\l{\cel{3}o di la elaboro di produkturo di qua onu volas uniformigar}
\l{\cel{3}la uzo o certigar la inter-idanteso perfekta, do la inter-}
\l{\cel{3}remplasableso sen desfacileso. III. Tipo legala di la}
\l{\cel{3}mezurili e ponderili lego-permisata. - DEFIRS.}
}
